\documentclass{article}
\usepackage{graphicx}
\usepackage{subcaption}
\usepackage[T2A]{fontenc}
\usepackage[utf8x]{inputenc}
\usepackage[english, ukrainian]{babel}
\usepackage{url}
\usepackage{latexsym,amsmath,amsfonts,amssymb,amsthm,mathrsfs}
\usepackage{cite,hhline}
\usepackage{hyperref}
\usepackage{graphicx,mwe}
\usepackage{wrapfig}
\usepackage{array}
\usepackage{tabularx, multirow}
\usepackage{longtable}
\usepackage{lipsum} % just for dummy text- not needed for a longtable
%\usepackage{minted}

\addto\captionsukrainian{\renewcommand{\figurename}{Мал.}}
\usepackage{book}					% стильовий файл журналу 

\title{Мій перший документ}
\date{2022-11-14}
\author{МАН }



\begin{document}
	\maketitle
	\newpage
	
\Title{Оформлення шаблону книги}\label{radap1xxx:FirstPage}
\Authors{Мосійчук В.С.}
\aff{Київський політехнічний інститут імені Ігоря Сікорського, м. Київ, Україна}
\Address{}{\href{mailto:mvs@ros.kpi.ua}{mvs@ros.kpi.ua}}
\AbsKeywords{У роботі представлені вимоги щодо оформлення статей для подання у збірник наукових праць ``Вісник Національного технічного університету України <<Київський політехнічний інститут>>. Серія Радіотехніка. Радіоапаратобудування''. Показано, що дотримання встановлених правил дозволить покращити вашу статтю .
}{Ключові слова}{правила оформлення, радіотехніка, радіоапаратобудування, \LaTeX}
	


\input{part1}

\input{part2}

\input{conclusion}

\input{lozko_part1}

\input{miziuk_part}

\input{Mylyar_part1}

%\input{javafx/chapter}

\section{Мова програмування С++}
\Authors{\textbf{Овсянніков Олександр Сергійович}}
\Address{}{\href{mailto:vinniaso@gmail.com}{vinniaso@gmail.com}}

\begin{wrapfigure}{L}{0.25\textwidth}
	\includegraphics[width=3cm]{ISO_C++_Logo.svg.png}\\
	\caption{\label{fig:frog1} Логотип.}
\end{wrapfigure} 
\textbf{Мова програмування С++} представляє собою високорівневу мову програмування загального призначення зі статичною типізацією, що підходить для створення найрізноманітніших програмних продуктів. На сьогоднішній день С++ є однією з найпопулярніших (див. \href{https://www.tiobe.com/tiobe-index/}{TIOBE Index}) та найпоширеніших мов програмування.
Мова C++ є еволюційним розвитком мови програмування C. Мова C була розроблена на початку 1970-х років в Bell Labs під час розробки операційної системи Unix. Автором мови C вважається \href{https://uk.wikipedia.org/wiki/%D0%94%D0%B5%D0%BD%D0%BD%D1%96%D1%81_%D0%A0%D1%96%D1%82%D1%87%D1%96}{Денніс Рітчі} (\textit{Dennis Ritchie}), якому не вистачало доступних на той момент інструментів мов програмування для написання коду операційної системи Unix. На початку 1980-х років датський програміст \href{https://uk.wikipedia.org/wiki/%D0%91%27%D1%8F%D1%80%D0%BD_%D0%A1%D1%82%D1%80%D0%B0%D1%83%D1%81%D1%82%D1%80%D1%83%D0%BF}{Б’ярн Страуструп} (\textit{Bjarne Stroustrup}), який на той час працював у компанії Bell Labs, розробив мову С++ як розширення до мови програмування С. Фактично спочатку мова C++ просто доповнювала мову С деякими можливостями об'єктно-орієнтованого програмування. І тому, сам Страуструп спочатку називав мову – «С з класами» (\textit{«C with classes»}). Згодом нова мова програмування стала набирати популярності. До неї були додані нові можливості, які робили її не просто доповненням до мови С, а й взагалі новою мовою програмування. У результаті в 1983 році мову програмування «С з класами» було перейменовано на С++. І з тих пір обидві мови програмування стали розвиватися незалежно одна від одної.

\begin{center}
	\textbf{Історія версій}
\end{center}
\begin{longtable}{| p{.1\textwidth} | p{0.8\textwidth} |} 
	\hline
	1985 & випущена перша комерційна версія мови С++, а також перше видання книги «Мова програмування C++», яка представляла перший опис цієї мови за відсутності офіційного стандарту \\ \hline 
	1989 & випущена нова версія мови C++ 2.0, що включала низку нових можливостей \\ \hline
	1998 & була зроблена перша спроба стандартизації цієї мови програмування організацією \href{https://uk.wikipedia.org/wiki/%D0%9C%D1%96%D0%B6%D0%BD%D0%B0%D1%80%D0%BE%D0%B4%D0%BD%D0%B0_%D0%BE%D1%80%D0%B3%D0%B0%D0%BD%D1%96%D0%B7%D0%B0%D1%86%D1%96%D1%8F_%D0%B7%D1%96_%D1%81%D1%82%D0%B0%D0%BD%D0%B4%D0%B0%D1%80%D1%82%D0%B8%D0%B7%D0%B0%D1%86%D1%96%D1%97}{International Organiztion for Standartization} (ISO). Перший стандарт отримав назву ISO/IEC 14882:1998 або скорочено C++98 \\ \hline
	2003 & було видано нову версію стандарту C++03 \\ \hline
	2011 & видано новий стандарт C++11, який містив безліч додавань та збагачував мову програмування С++ великою кількістю нових функціональних можливостей \\ \hline
	2014 & було випущено невелике додавання до стандарту, відоме також, як C++14 \\ \hline
	2017 & cтандарт С++17 був офіційно опублікований. Було внесено ряд змін в мову програмування, зокрема додано декілька нових бібліотек, алгоритмів  та класів до стандартної бібліотеки шаблонів (STL), поліпшено роботу з лямбда-виразами, оновлено бібліотеку <thread> для мультипотокового програмування \\ \hline
	2020 & cтандарт C++20. Ця версія має набагато більше нововведень в порівнянні з C++14 та C++17. Зокрема, ця версія вводить «концепції», «співпрограми», новий оператор порівняння <=>, тощо \\ \hline
\end{longtable}

\textbf{С++} є потужною мовою програмування, яка успадкувала від мови С багато можливостей по роботі з пам'яттю. Тому нерідко мова С++ знаходить своє застосування у системному програмуванні, зокрема, при створенні операційних систем, драйверів, різних утиліт, антивірусів тощо. До речі, ОС Windows переважно написана на мові програмування С++. Але лише системним програмуванням застосування цієї мови не обмежується. 

\textbf{С++} можна використовувати у програмах будь-якого рівня, де важливі швидкість роботи та продуктивність. Нерідко вона застосовується для створення графічних додатків, різних прикладних програм. Також особливо часто її використовують для створення ігор із насиченою візуалізацією. Крім того, мова С++ застосовується в мобільному програмування, у веб-розробці для створення веб-додатків або будь-яких допоміжних сервісів, які обслуговують веб-додатки. Загалом С++ – мова програмування широкого користування, якою можна створювати практично будь-які види програмного забезпечення.



\section
\title\centering\textbf{Обухівка}

\date{6 грудня 2022 рік}

\author\textbf{Старотіторов Валерій Петрович}

\maketitle
%\begin{document}
	Обухівка розташована в центральній частині Дніпропетровської області на лівому березі Дніпровського водосховища на Дніпрі. Неподалік Обухівки знаходиться місце впадіння нового русла річки Оріль у Дніпро. Зі сходу до Обухівки впритул прилягають лівобережні околиці міста Дніпра. На східній околиці з озера Чередниченкове витікає річка Шпакова. На півночі межує із селом Горянівське.

Обухівка розташована у фізико-географічній зоні Придніпровська низовина. Висота над рівнем моря у селищі — 50-55 метрів.

Довкола селища розташовані Обухівські озера.

На захід від Обухівки лежать землі Дніпровсько-Орільського природного заповідника.
\begin{figure}
	\centering
	\includegraphics[width=0.7\linewidth]{25}
	\caption{Арт-об'єкт "Я кохаю Обухівку"}
	\label{fig:25}
\end{figure}
\begin{center}
	\begin{tabular}{|c | c|} 
		\hline
		Засноване & II п. XVIII ст.\\ [0.5ex] 
		\hline\hline
		Площа & 14,9 км² \\ 
		\hline
		Населення & 8266 \\
		\hline
		Густота & 554 осіб/км²\\
		\hline
		Висота над рівнем моря & 54 м \\
		\hline
		Телефонний код & +380 56(2) \\ [1ex] 
		\hline
	\end{tabular}
\end{center}
\url{https://uk.wikipedia.org/wiki/%D0%9E%D0%B1%D1%83%D1%85%D1%96%D0%B2%D0%BA%D0%B0_(%D1%81%D0%BC%D1%82)}

Історія
На території Обухівки досліджено кілька поселень доби бронзи, катакомбної та зрубної культур (кінець ІІІ — початок І тисячоліття до н. е.).

Засноване в другій половині XVIII століття. Першою назвою поселення була Обухівка, яка, за неофіційними даними, походить від прізвища першого поселенця — козака Андрія Обуха.

Після утворення Катеринославської губернії Обухівка ввійшла до складу Кам'янської волості Новомосковського повіту. У 1886 році тут мешкало 1930 осіб з 351 господарським подвір'ям, була церква, 3 магазини.

В перші роки радянської влади в Обухівці, як і у всій країні, пройшла колективізація. Станом на 1925 рік в селі мешкало 5 246 жителів, проводилось два ярмарки, було дві школи, хата-читальня[1]. Обухівка входила до складу Кам'янського району Катеринославської округи.

У 1938 році, за рішенням загальних зборів селян Обухівки, вирішено перейменувати село на Кіровське на честь радянського державного діяча Сергія Кірова.[2] Згодом до селища була приєднано придніпровське село Сугаківка. Обухівка увійшла до складу новоутвореного Дніпропетровського району.

Після Другої світової війни Кіровське отримало статус селища міського типу. Тут була розташована центральна садиба Дніпропетровської птахофабрики. У другій половині 20 століття у Кіровському побудовано понад 1 200 нових житлових будинків, Палац культури, дві бібліотеки, лікарська амбулаторія.

Деякий час Обухівка підпорядковувалася Дніпропетровській міській раді, з 1965 року, після відновлення Дніпропетровського району, знову входить до нього.

% TODO: \usepackage{graphicx} required
\begin{figure}[h]
	\centering
	\includegraphics[width=0.8\linewidth]{23}
	\caption{\label{fig:frog3}Храм Святих мучениць Віри, Надії, Любові та матері їх Софії}
	\label{fig:23}
\end{figure}

Сучасність
Більшість (90 %) працюючого населення Обухівки працює у місті Дніпрі.

У селищі є дві середні загальноосвітні школи, дві амбулаторії загальної практики—сімейної медицини, будинок культури.

Дитячий садок, збудований у 2018 році, у лютому 2020 року номінований організаторами конкурсу Mies van der Rohe Award на премію ЄС у галузі сучасної архітектури.[3]
\begin{figure}
	\centering
	\includegraphics[width=0.8\linewidth]{24}
	\caption{	\centering Обухівський дитсадок Дізнатись більше: https://litsa.com.ua/dityachij-sadok-v-obuxivci-nominuvali-na-premiyu-yes/}
	\label{fig:24}
\end{figure}

Обухівка відіграє роль «дачного селища» та місця відпочинку для жителів міста Дніпра, оскільки за часів радянської влади тут було розміщено багато дач та баз відпочинку (в основному для оздоровлення робітників місцевих заводів та фабрик). У 2017 році в Обухівці створено громадську організацію "Дитячий Туристичний Клуб «БАФ», якій об'єднав понад 100 активних громадян селища. У 2018 році в Обухівці створено «Молодіжний центр А. Т.О. М.», центр є майданчиком для зустрічей активної молоді селища

Герб
Орлан і пернач засвідчують давні козацькі традиції, пов’язують зі символікою Орільської полкової паланки, на території громади річка Оріль впадає у Дніпро. Орлан є символом волі й благородства. Три зірки означають сходження тут меж колишніх Орільської, Протовчанської та Кодацької паланок, вказують близькість до міста Дніпра (на гербі якого такі ж зірки) та характеризують об’єднану громаду, до складу якої входять три населені пункти. Хвиляста основа підкреслює розташування поселення біля річки Дніпра. Зелене поле символізує щедру природу, Дніпровсько-Орільський природний заповідник. Дубова та ялинова гілки підкреслюють багатство місцевої флори, є символом міцності та сили.

Транспорт
Через селище проходять автомобільні дороги Т 0404 і Н31.

Обухівка сполучена приміським транспортом (маршрутками) з Дніпром.

Пам'ятки
Храм на честь Воздвиження Чесного і Животворящого Хреста Господнього
Поблизу Обухівки наявний регіональний ландшафтний парк Дніпровські ліси.
Відомі люди
Уродженцями села є

Варфоломєєв Володимир Олександрович (1984—2014) — старший солдат Збройних сил України, учасник російсько-української війни.
Лисенко Олег Ігорович (1989—2016) — молодший сержант Збройних сил України, учасник російсько-української війни.
	
\end{document}



\end{document}
